\section{Experimental Setup}
\label{sec:setup}

\paragraph{Testbed.} We use GENIUS~\cite{kim2025genius} as testbed throughout: a CLIP-SF (CLIP with score fusion) encoder produces joint image/text embeddings, an 8-level RQ (plus a modality-indicator level, codebook size 4096) quantizes each into a semantic ID, and a T5 decoder generates a query's target ID via constrained beam search (50 beams) over a trie of valid candidate IDs, so decoding can never emit a nonexistent ID.

\paragraph{Datasets as contrasting regimes.} We train and evaluate on MSCOCO~\cite{lin2014mscoco}, VisualNews~\cite{liu2021visualnews}, and FashionIQ~\cite{wu2021fashioniq} independently. MSCOCO and VisualNews are both broad-domain, bidirectional captioning retrieval (text$\to$image, image$\to$text) drawn from different image/text distributions (everyday photography vs.\ news photography and captions); FashionIQ is composed retrieval (image + modification text $\to$ target image) in one narrow domain, its pool dense with visually similar garments. Candidate pools differ substantially in size -- MSCOCO $713{,}679$, VisualNews $199{,}903$ ($100$K image, $99.9$K text), FashionIQ $74{,}380$ -- so we never compare absolute Recall across datasets, only relative effects within each. Pool size is also not what drives the dataset-dependence we report: re-running MSCOCO's ID-structure diagnostics on a subsample size-matched to FashionIQ's pool ($74{,}380$ candidates) still yields $2.1$--$4.0\times$ FashionIQ's codebook utilization at every $\lambda$, with far steeper collision growth. For VisualNews, no official M-BEIR held-out test partition exists in our downloaded data, so we repurpose the union pool's validation split (never seen during training) as the test set -- a disclosed deviation from the official-test-split convention used for MSCOCO and FashionIQ.

\paragraph{Pipeline and controls.} For each dataset we fix the CLIP-SF encoder, extract candidate embeddings once, then train RQ tokenizers differing only in balance-regularization strength $\lambda$ -- four points $\{0,0.3,1.0,3.0\}$ (vanilla/weak/medium/strong) for MSCOCO and FashionIQ, two points $\{0,3.0\}$ for VisualNews -- holding architecture, level count, codebook size, decoder, and decoding settings fixed. Each tokenizer's IDs retrain the same decoder from scratch. Headline Recall numbers are means over multiple decoder-training seeds per variant (5 for MSCOCO vanilla/strong, 3 elsewhere); MSCOCO's headline gain is further verified across 3 independently retrained RQ tokenizers (Section~\ref{sec:finding1}), not just decoder seeds.

\paragraph{Scope: relative effects under a fixed, non-SOTA-reproducing pipeline.} Because the officially released GENIUS decoder is not public, we pair the officially released Stage-1 RQ checkpoint with our own decoder trained on the full $\sim$1.3M-query, 16-task M-BEIR union, and compare against the published paper's numbers: FashionIQ lands above the published figure ($18.3$--$18.6\%$ R@10 vs.\ $13.1\%$, no reranking) while MSCOCO T$\to$I R@1 lands far below it ($8.1\%$ vs.\ $40.1\%$). We do not read the FashionIQ result as proof the pipeline is faithful -- overshooting on one dataset and undershooting on another is itself a reason for caution -- only as evidence that the gap is not a uniform implementation failure. The MSCOCO shortfall is consistent with capacity dilution across 16 joint tasks~\cite{pradeep2023scale}, plus our lack of validation-based checkpoint selection and hard negatives -- training-recipe gaps we disclose rather than close further. This union-trained reproduction is a separate regime from every $\lambda$-sweep comparison in this paper: each sweep variant instead trains a \emph{fresh} RQ tokenizer and decoder from scratch, per dataset, with $\lambda$ as the only difference, and is never exposed to the union-training gap above. We therefore report the effect of regularization \emph{relative to a matched vanilla baseline under an identical from-scratch pipeline}, not GENIUS's absolute state-of-the-art performance, and we do not pursue closing the reproduction gap further within this study.
